\chapter{Discussion}
\label{chap:discussion}

This thesis aimed to clarify the assumptions and the finite-sample performance of two 
approaches to address 
time-constant unmeasured confounding in the context of time-varying treatments. 
Although the simulation results show improvements in the bias of the 
causal effect estimates for both methods, there are still some limitations and 
open questions which make these methods not yet ready for 
widespread use in practice. The fundamental requirement behind both of these methods 
is another form of ignorability assumption which includes a time-invariant latent variable 
in the conditioning set which make the treatment assignment independent of the potential 
outcomes. 

In FE-MSM this requirement is directly assumed and is formalized by 
Assumption~\ref{ass:FE-unit-specific-ignor}. The restriction they put in this assumption 
on the unmeasured confounder is that using a latent scalar variable $\alpha_i$, 
all the 
time-invariant unmeasured confounding in the treatment process 
can be captured. Effectively, this assumption restricts the class 
of DGPs which this method can address to the class of DGPs where the unmeasured confounder 
only enters the treatment process as a scalar or, at least, where those confounders can be summarized 
into one scalar. This is a strong structural assumption which cannot be verified 
using observed data. For instance, this method cannot be used in cases where 
the unmeasured confounding modifies not only the intercept but also 
the effect of other covariates or treatments (interaction effects). 

On the other hand, in SeqGPLVM, this requirement 
is justified by means of assuming a set of auxiliary 
assumptions (Assumptions~\ref{ass:SG-Marcovianity}-\ref{ass:SG-Consistency}). 
Among these assumptions, the assumption of consistency of the latent variable 
plays the most crucial role. This assumption is not only impossible to verify in 
real-world problems, but it is not straightforward to establish even in simulation studies where 
the DGP is known. To verify this assumption for a given DGP, one would need to show that the 
relevant latent confounder is uniquely identified, up to trivial transformations, 
by the observed data. Note that the notion of identifiability considered here 
is closer to structural identifiability, since the issue is whether the 
latent confounder is uniquely implied by the data-generating structure, 
rather than whether a parameter can be estimated within a fixed statistical model. 
This is a much 
stronger requirement than merely showing that there exists some latent-variable representation 
of the treatment process. In general, different latent structures may induce the same observed 
distribution of the treatments, and therefore the same observed data may be compatible with 
multiple candidate substitute confounders. In such a case, even if a latent-variable model fits 
the treatment process well, it does not follow that the recovered latent variable corresponds to 
the confounder needed for causal identification. In other words, SeqGPLVM 
may recover a latent factor that explains the dependence structure among the treatments, 
but this alone does not guarantee that this latent factor is the correct one for the 
causal problem. This is a fundamental issue for all methods that rely on latent 
variable models for deconfounding, and it is not specific to SeqGPLVM.

Consequently, the favorable simulation results for SeqGPLVM should be 
interpreted with caution. This method is shown to work in a specific DGP and for larger 
sample sizes, but the satisfaction of the concistency assumption is left for 
future work. More broadly, the method shifts the 
burden of the usual sequential ignorability assumption to a different and 
arguably less transparent set of 
structural assumptions about the existence, uniqueness, and causal relevance of a 
latent substitute 
confounder. Despite these theoretical limitations, the flexibility of SeqGPLVM 
in modeling complex relationships 
between the observed covariates, treatments, and the latent confounder for longitudinal data, 
can be studied further as a non-parametric latent variable model for time series data. 

A future direction is to add dynamics to the latent variable structure. Also another 
interesting direction is to implement an amortized version of SeqGPLVM which can 
be more efficient for new data by parametrizing the posterior distribution over the 
latent variables using a neural network.

In conclusion this thesis provided a common framework for understanding the assumptions 
and the finite-sample performance of two methods for addressing time-constant unmeasured 
confounding in time-varying treatments settings. 
It highlighted the assumptions and limitations of these methods and 
the need for further theoretical and empirical work to understand their 
properties and applicability. The results of this thesis suggested 
evidence that both FE-MSM and SeqGPLVM perform similarly for longer treatment sequences.  The 
results of this thesis provided emperical evidence that both FE-MSM 
and SeqGPLVM can be useful tools for addressing time-constant unmeasured confounding in 
certain settings to reduce bias and improve coverage compared to standard methods 
that do not account for unmeasured confounding. A promising future research 
direction is to 
investigate the theoretical properties of the fixed effects estimates by FE-MSM and the 
latent variables estimated by SeqGPLVM to understand the nature of these estimates 
and their relationship to the true unmeasured confounder. This could be pursued by 
further exploring the performance and the estimation of the latent variables 
and the fixed effects 
in simulation studies using alternative DGPs 
that impose a richer structure on the unmeasured confounding. 
For example, one could consider 
a hierarchical DGP in which individuals are first assigned to latent 
classes and the unmeasured confounder 
is then generated conditional on class membership. Such a design would 
induce structured heterogeneity in 
the population and would make it possible to investigate whether the fixed effects 
in FE-MSM and the latent 
variables inferred by SeqGPLVM recover only coarse subgroup structure or a 
quantitatively meaningful 
approximation to the underlying confounder. 
In this sense, the simulation framework suggested by 
\citeA{bouchattaouiCausalDynamicVariational2025} may provide a useful starting point, 
even though that 
paper focuses on unobserved adjustment variables rather than unobserved confounders.

Another idea inspired by \citeA{bouchattaouiCausalDynamicVariational2025} is to investigate 
whether the consistency assumption for SeqGPLVM can be imposed also in practice by imposing 
a more concentrated,
approximately deterministic posterior over the latent variables. 
The theory of deconfounding requires 
that the latent substitute confounder is a deterministic function of the 
treatment sequence 
while in practice the posterior distribution over the latent variables is 
typically diffuse, especially for smaller sample sizes. If successful, 
such an approach could lead 
to more stable latent representations 
and reduce the sensitivity of the resulting causal effect estimates to 
posterior uncertainty in the inferred confounder.

\begin{comment}
\textbf{DISCUSSION}
\hl{While this MI-style propagation provides a principled empirical uncertainty summary,
it does not close the main theoretical gap in the sequential deconfounding approach.
First, it does not by itself deliver a large-sample result for the \emph{overall}
two-stage estimator---i.e., conditions under which the plug-in (or draw-averaged)
MSM estimator is consistent and $\sqrt{N}$-asymptotically normal with a valid variance
formula---of the kind established for FE--IPTW by Blackwell and Yamauchi.
Second, because the outcome stage is not modeled jointly with the treatment model,
the distribution of $\{\hat{\gamma}^{(s)}\}$ should be interpreted as uncertainty
propagation rather than a coherent Bayesian posterior for $\gamma$; consequently,
interval calibration is not automatically guaranteed from Bayesian first principles.
Third, even perfect variance combination does not address potential sources of
systematic error---such as misspecification of the latent treatment model/CMM,
violations of positivity leading to unstable weights, or approximation error from
variational inference---which can affect both bias and coverage.}


\begin{itemize}
    \item discuss about the vocablurary choice of "unmeasured" vs "hidden"
    \item testing against misspecified latent space dimension (I think bica does this?)
    \item implications of having one GP per time step vs one big GP for all time steps (maybe cite 
    some GP time series work?)
    \item discuss the none parametric nature of GP and its implications, specifically in the context of IPTW and the assumption of correct propensity model specification
    \item discuss about the amortized version of GPLVM (perhaps in the appendix)
    \item \hl{what} \cite{hattSequentialDeconfoundingCausal2024} \hl{did seems very wrong. it seems that they extract the latent variables from the treatment model and then use them as if they were observed covariates in the MSM. The worrisome part is that they are using IPTW which means after extracting the esimated latent substitute, they fit another propensity model incluing the estimated treatments so basically on the left hand side we have $a_it$ and the on the right hand side we have a function of $a_it$. This seems very wrong. We should discuss this. What I did, which seems more correct, is to use the estimated propensities from the seqgplvm treatment model directly in the MSM. Althoughm, this can also cause some problems because seqgplvm is a very powerful model and it can overfit the treatment assignment and lead to extreme weights.}
    \item another break point of this whole story of deconfounding with latent variables is the assumption of pinpointability. in the context of the time varying treatments, this assumption means that the latent variables estimated should stay the same for all hypothetical treatment assignments. I need to read D'amor critique of deconfounding paper again. 
    \item another big issue is that, their second propensity model is the msm stage is not even correctly specified! Although they removed the explicit equations they used for the MSM stage, from their previous version, I remeber that they used a model that was not the same as the DGP. 
    \item Another point is that, the selling point of Hatt et al was that SeqGPLVM is flexible in terms of 
    captuing non linear relationships between the latent confounders, observed covariates and treatment assignments. However, 
    because of the pinpointability assumption, there is a high chance that the those more complicated relationships in the dgp, 
    cannot satisfied the pinpointability assumption. It would be nice if I can find a dgps that is not very complicated but 
    violates the pinpointability assumption. Maybe reading D'amor critique paper again can help with this.
    \item discuss about the SeqGPLVM architecture choice of having separate parameters for each time step vs having shared parameters across time steps and its
    implications. Note that this is not the most efficient way to parameterize a time series model if we believe that the dgp's paramters do not change over time.  
    \item discuss about the big assumption \cite{blackwellEffectPoliticalAdvertising2024} 
    make about $Z_i = f(U_i)$ and perhaps what effectively \cite{hattSequentialDeconfoundingCausal2024} 
    is doing to not assume this but instead in their theorwm 1, they are tyring to show that 
    their $Z_i$, is enough to retain the sequential ignorability assumption. But a future 
    direction is to investigate if one results in another in the sense that if 
    $Z_i$ is enough to retain the sequential ignorability assumption, then it must be the case that $Z_i = f(U_i)$ for some function $f$.
    \item advantage using seqgplvm is that it let the uncertainty in the 
    estimation of propensity scores to be propagated to the final 
    estimation of the causal effect. 

    \item in this work we show via simulations that SeqGPLVM can result in similar results as 
    FE-MSM and this motivates to work on the consistency theory of SeqGPLVM similar to 
    what \cite{blackwellEffectPoliticalAdvertising2024} did for FE-MSM. more specifically, 
    we need to show that the propensity scores estimated by SeqGPLVM are 
    consistent estimates of the true propensity scores. 

    \item discuss that it would be better if we use the amortized version of 
    SeqGPLVM since now for new data we need to refit only the part of the model 
    corresponding to the latent variable while freezing the rest of the parameters. 
    but via amortized inference, we can directly get the latent variable for new data without refitting any part of the model.
\end{itemize}

\section{Problem of reusing the estimated latent variables in a second stage outcome model involving propensity score estimation}
\label{sec:hatt-problems}
As mentioned in section \ref{sec:seqgplvm}, in \cite{hattSequentialDeconfoundingCausal2024}, the authors treat 
SeqGPLVM as a factor model whose purpose is to estimate the posterior over the latent confounder substitutes, 
\(p(\mathbf{z}_i \mid \bar{a}_{iT}, \bar{\mathbf{x}}_{iT})\). Then, the authors proceed to use samples 
from this posterior distribution as if they were observed covariates in a second stage outcome model. 
Specifically, they suggest fitting two propensity score based models: 
marginal structural model (MSM)\cite{robinsEstimationCausalEffects2008} and 
Recurrent Marginal Structural Networks \cite{limForecastingTreatmentResponses2018}. The 
problem with this approach is that in both of these models, the authors fit a second propensity score model 
that includes the estimated confounder substitutes \(\mathbf{z}_i\) as covariates. However, since \(\mathbf{z}_i\) 
are estimated from the treatment assignments \(\bar{a}_{iT}\), including them as covariates in a second propensity score 
model that aims to predict \(\bar{a}_{iT}\) creates a circular dependency. More concretely, 
in the second stage propensity score model, we are trying to estimate 
\(p(A_{it} \mid \bar{A}_{i,t-1}, \bar{\mathbf{X}}_{it}, \mathbf{z}_i)\) where \(\mathbf{z}_i\) is a function of 
\(\bar{a}_{iT}\) and \(\bar{x}_{iT}\). 
This circular dependency can lead to biased estimates of
using inverse probability of treatment weighting (IPTW), where the propensity scores are estimated by including the estimated confounder substitutes $\mathbf{z}_i$ as covariates.

Also note that the correct way of using the posterior distribution over 
the latent confounder substitutes \(p(\mathbf{z}_i \mid \bar{a}_{iT}, \bar{\mathbf{x}}_{iT})\)
is to integrate over this distribution when estimating the causal effect, i.e.,
\begin{equation}
    \mathbb{E}\left[Y_{iT}(\bar{d}_{T-k:T})\right] = \int \mathbb{E}\left[Y_{iT}(\bar{d}_{T-k:T}) \mid \mathbf{z}_i\right] p(\mathbf{z}_i \mid \bar{a}_{iT}, \bar{\mathbf{x}}_{iT}) d\mathbf{z}_i
\end{equation}
which can be approximated using Monte Carlo integration by drawing samples from the posterior distribution, 
calculate the conditional propensity scores by the SeqGPLVM likelihood for the 
observed treatment, and then use these propensity scores in the IPTW estimation of the MSM.
Repeat this for multiple samples from the posterior and average the results to obtain the final estimate of the causal effect. 

Note that in \cite{bicaTimeSeriesDeconfounder2020} and \cite{wangBlessingsMultipleCauses2019} 
They also follow a similar two stage procedure where they first estimate the latent confounders
using a factor model and then use these estimated confounders as observed covariates in the outcome model. 
However, they do not use IPTW for estimating the causal effect in the outcome model. 
Instead, they use G-formula which does not involve fitting a second propensity score model.
\end{comment}